%!TEX root = ../../main.tex
% Figures for the paper, generated by scripts/48_latex_tables.py.
% Images are referenced as \tabfig/<file>; copy the PNG files next to this file.

\begin{figure}[htbp]
\centering
\subfloat[All cells]{\includegraphics[width=0.49\textwidth]{\tabfig/es_log_emp_all.png}}
\hfill
\subfloat[Cells aged 15--29]{\includegraphics[width=0.49\textwidth]{\tabfig/es_log_emp_young.png}}
\caption{Event-study coefficients, log cell employment, eleven countries pooled}
\label{fig:es_pooled}
\floatfoot{\footnotesize \textit{Notes:} Coefficients $\beta_k$ on 1[quarter $= k$] $\times$ High (top exposure quintile) with 95 percent confidence bands; reference 2022Q4 (dashed line). Cell fixed effects and country $\times$ age group $\times$ sex $\times$ quarter fixed effects; weights are mean 2022 cell employment; standard errors clustered by country $\times$ occupation. 2021 quarters cover the countries observed then (Brazil, Mexico, Argentina, Bolivia, Georgia and the Philippines from 2021Q1; Ecuador and Uruguay from 2021Q3); 2026 quarters cover the countries already released.}
\end{figure}

\begin{figure}[htbp]
\centering
\includegraphics[width=0.92\textwidth]{\tabfig/emp_index_q5_panel.png}
\caption{Employment by quintile of generative-AI exposure, 2022 $=$ 100, by country and pooled}
\label{fig:index}
\floatfoot{\footnotesize \textit{Notes:} Employment summed over all cells in the quintile, divided by its 2022 mean. Quintiles are employment-weighted at the 2022 baseline and defined at each country's cell digit level; quintile 5 is the most exposed. The pooled series is the employment-weighted mean of the country indices over the countries observed in each quarter. Bolivia and Georgia have few cells and noisy indices.}
\end{figure}

\begin{figure}[htbp]
\centering
\includegraphics[width=0.8\textwidth]{\tabfig/es_new_hire_share_all.png}
\caption{Event-study coefficients, new-hire share, all cells}
\label{fig:newhire}
\floatfoot{\footnotesize \textit{Notes:} Outcome: weighted share of employed persons in the cell who started their current job within the previous twelve months, in share points. Seven countries with a tenure question: Brazil, Colombia, Ecuador, Argentina, South Africa (all quarters), Uruguay (first interviews) and Mexico (first-interview quarters). Coefficients $\beta_k$ on 1[quarter $= k$] $\times$ High (top exposure quintile) with 95 percent confidence bands; reference 2022Q4 (dashed line). Cell fixed effects and country $\times$ age group $\times$ sex $\times$ quarter fixed effects; weights are mean 2022 cell employment; standard errors clustered by country $\times$ occupation. The quarter coefficients average to the difference-in-differences estimate in Table~\ref{tab:newhire}.}
\end{figure}

\begin{figure}[htbp]
\centering
\includegraphics[width=\textwidth]{\tabfig/es_new_hire_share_panel_country.png}
\caption{New-hire share event study by country, all cells}
\label{fig:newhire_country}
\floatfoot{\footnotesize \textit{Notes:} Each panel re-estimates the model within the country; 95 percent bands. Outcome and specification as in the pooled figure. South Africa has 46 clusters and Argentina 31; Mexico contributes first-interview quarters only.}
\end{figure}

\begin{figure}[htbp]
\centering
\includegraphics[width=0.9\textwidth]{\tabfig/itbpo_quarterly_portrait.png}
\caption{IT and business-process services, quarterly, countries with sufficient sector sample}
\label{fig:itbpo_quarterly}
\floatfoot{\footnotesize \textit{Notes:} Brazil, Colombia, the Philippines, India and Uruguay, and a pooled series over Brazil, Colombia, the Philippines and Uruguay for quarters where all four are observed (India is left out of the pool: its block is four times the others combined and its sample design changes in 2025). Dashed line at 2022Q4. India before 2025Q2 is first-visit records only.}
\end{figure}

\begin{figure}[htbp]
\centering
\includegraphics[width=\textwidth]{\tabfig/itbpo_occupation_split.png}
\caption{Occupation split within IT and business-process services, 2022 $=$ 100}
\label{fig:split}
\floatfoot{\footnotesize \textit{Notes:} Within the sector, employment in professional and technical occupations (ISCO major groups 1 to 3) and in clerical and service occupations (major groups 4 and 5), indexed to the 2022 annual average, against all employment. 2026 covers one or two quarters.}
\end{figure}

\begin{figure}[htbp]
\centering
\subfloat[All cells]{\includegraphics[width=\textwidth]{\tabfig/es_log_emp_panel_country.png}}
\\
\subfloat[Cells aged 15--29]{\includegraphics[width=\textwidth]{\tabfig/es_log_emp_panel_country_young.png}}
\caption{Event-study coefficients, log cell employment, by country}
\label{fig:es_country}
\floatfoot{\footnotesize \textit{Notes:} Each panel re-estimates the model within the country; 95 percent bands. Coefficients $\beta_k$ on 1[quarter $= k$] $\times$ High (top exposure quintile) with 95 percent confidence bands; reference 2022Q4 (dashed line). Cell fixed effects and country $\times$ age group $\times$ sex $\times$ quarter fixed effects; weights are mean 2022 cell employment; standard errors clustered by country $\times$ occupation.}
\end{figure}

\begin{figure}[htbp]
\centering
\subfloat[All cells]{\includegraphics[width=0.49\textwidth]{\tabfig/es_log_emp_IND_to2024Q4_all.png}}
\hfill
\subfloat[Cells aged 15--29]{\includegraphics[width=0.49\textwidth]{\tabfig/es_log_emp_IND_to2024Q4_young.png}}
\caption{India (annex): event-study coefficients, log cell employment, urban sample, 2021Q3 to 2024Q4}
\label{fig:india}
\floatfoot{\footnotesize \textit{Notes:} PLFS urban first-visit records, three-digit cells; the series stops before the 2025 redesign. Coefficients $\beta_k$ on 1[quarter $= k$] $\times$ High (top exposure quintile) with 95 percent confidence bands; reference 2022Q4 (dashed line). Cell fixed effects and country $\times$ age group $\times$ sex $\times$ quarter fixed effects; weights are mean 2022 cell employment; standard errors clustered by country $\times$ occupation.}
\end{figure}
